\documentclass[10pt,a4paper]{amsart}
\usepackage[margin=19mm]{geometry}
\usepackage{amsmath,amssymb,mathtools}
\usepackage[scr=rsfs,cal=euler]{mathalfa}
\usepackage{xfrac}
\usepackage[unicode,hidelinks,bookmarks=false]{hyperref}
\newcommand{\ie}{i.e.\ }
\newcommand{\cf}{cf.\ }
\newcommand{\resp}{respectively\ }
\newcommand{\Subsection}{\subsection}
\providecommand{\nobreakdash}{\nobreak\hbox{-}}
\setlength{\emergencystretch}{2em}
\multlinegap0pt
\usepackage{amssymb}
\usepackage{xcolor}
\usepackage{bm}
\usepackage{algorithm}\usepackage{algpseudocode}
\usepackage[shortlabels]{enumitem}
\usepackage{mathtools}
\usepackage{tcolorbox}
\usepackage{dsfont}
\usepackage{tikz}
\newtheorem{lemma}{Lemma}
\newcommand{\Tnorm}[1]{\left\vert\mkern-3mu\left\vert\mkern-3mu\left\vert #1 \right\vert\mkern-3mu\right\vert\mkern-3mu\right\vert}
\newcommand{\bH}{{\bf H}}
\newcommand{\bI}{{\bf I}}
\newcommand{\bSigma}{\bm{\Sigma}}
\newcommand{\bbR}{{\mathbb R}}
\newcommand{\bbZ}{{\mathbb Z}}
\newcommand{\be}{\begin{equation}}
\newcommand{\bean}{\setlength\arraycolsep{1pt}\begin{eqnarray*}}
\newcommand{\bfe}{{\bf e}}
\newcommand{\bx}{{\bf x}}
\newcommand{\ee}{\end{equation}}
\newcommand{\eean}{\end{eqnarray*}}
\newcommand{\tnorm}[1]{\vert\mkern-3mu\vert\mkern-3mu\vert #1 \vert\mkern-3mu\vert\mkern-3mu\vert}
\title{\vspace*{-9mm} Nonparametric undirected graphical model selection \\ using diffusion models}
\author{Hyeok Kyu Kwon, Myeonggu Kang, Minwoo Chae, Wanjie Wang}
\date{Source: arXiv:2606.08468v1 (CC BY 4.0)}
\begin{document}
\maketitle

\noindent\emph{Source and licence.} \vspace*{-9mm} Nonparametric undirected graphical model selection \\ using diffusion models. Version arXiv:2606.08468v1 (CC BY 4.0), distributed by arXiv under the Creative Commons Attribution 4.0 International licence, http://creativecommons.org/licenses/by/4.0/, as recorded in the arXiv licence metadata for this exact version and shown on the abstract page https://arxiv.org/abs/2606.08468v1. Full licence notice: \texttt{licenses/arxiv-2606.08468-CC-BY-4.0.txt}.\par\smallskip

\noindent\textit{Editorial scope.} Counted unit: the article's lemma lem1 (source0.tex lines 1495--1505). The article's complete original proof of it is delivered with it (source0.tex lines 1507--1583, one \texttt{proof} environment of 77 lines). No mathematics is written, repaired or re-derived: the statement and the proof are the article's own lines, in the article's own notation, and no retained line is substituted or re-flowed.  No further source-local context is needed: every cross-reference of every retained line resolves inside the two delivered spans.  Nothing was omitted from any retained span, so the package carries no omission record.  No retained line contains a citation, so the wrapper carries no bibliography and no bibliography entry.  No retained line contains a figure environment, an image-inclusion command or drawing code, so no image is delivered and the package declares no asset dependency.  Editorial formatting only, in addition: an AMS \texttt{amsart} preamble that restates the article's own declarations and macros used by the retained lines, namely the environments \texttt{lemma} on separate counters, together with the standard packages those lines need.  The retained environments print in this document as Lemma 1 in order of appearance, whereas the article prints the same items with the numbering of its own full document.  The complete native source of the article as distributed in the arXiv e-print for this exact version is bundled unchanged as \texttt{sources/202334/source0.tex}.
\begin{lemma}\label{lem1}
Let $K > 0$ be given, and consider an undirected graph $G$ with vertex set $\{1,\ldots,D\}$.
Suppose that a symmetric matrix $\bH  \in \bbR^{D \times D}$ satisfies $\vert H_{ij} \vert \leq K$ for all $i,j \in [D]$ and $H_{ij} = 0$ when $d_{G}(i,j) > 1$.
Then, for every $0 \leq \epsilon < (DK)^{-1}$,
\bean
\left\vert \left\{ \left( \bI_{D} - \epsilon \bH \right)^{-1} \right\}_{ij} - \epsilon^{d_{G}(i,j)} \left( \bH^{d_{G}(i,j)} \right)_{ij} \right\vert 
\leq \frac{ (\epsilon D K)^{d_{G}(i,j)+1} }{1-\epsilon D K}
, \quad \forall i, j \in [D] \text{ with } d_{G}(i,j) < \infty,
\eean
and $\{ (\bI_{D} - \epsilon \bH)^{-1}\}_{ij} = 0$ when $d_{G}(i,j) = \infty$.
\end{lemma}

\begin{proof}
By the Cauchy–Schwarz inequality, we have $(\sum_{i=1}^{D} \vert x_i \vert)^2 \leq D \| \bx \|_2^2$ for any $\bx = (x_1,\ldots,x_{D}) \in \bbR^{D}.$
A simple calculation yields that
\bean
\bx^{\top} \bH \bx  = \sum_{i,j \in [D]} H_{ij} x_{i} x_{j}
\leq K  \sum_{i,j \in [D]} \vert x_i \vert \vert x_j \vert = K \| \bx \|_1^2 \leq D K \| \bx \|_2^2.
\eean
It follows that, for any $\epsilon > 0$,
\bean
\bx^{\top} \left( \bI_{D} - \epsilon \bH \right) \bx 
= \| \bx \|_2^2 - \epsilon \left( \sum_{i,j \in [D]} H_{ij} x_i x_j \right)
\geq (1-\epsilon D K)  \| \bx \|_2^2.
\eean
Hence, for every $0 \leq \epsilon < (DK)^{-1}$, the matrix $\bI_{D}-\epsilon \bH$ is positive-definite.
Let
\bean
\bSigma
=
\left( \bI_{D} - \epsilon \bH \right)^{-1}.
\eean
Then, for every $N \in \bbZ_{\geq 0}$, a simple calculation yields that
\bean
\bSigma - \sum_{k=0}^{N} \epsilon^{k} \bH^{k} = \epsilon^{N+1} \bSigma \bH^{N+1}
\eean
because $ \bSigma^{-1}  (\sum_{k=0}^{N} \epsilon^{k} \bH^{k} ) = ( \bI_{D} - \epsilon \bH ) (\sum_{k=0}^{N} \epsilon^{k} \bH^{k} )
= \bI_{D} - \epsilon^{N+1} \bH^{N+1}.$
For any $i,j \in [D]$,
\bean
\left\vert \left( \bSigma \bH^{N+1} \right)_{ij} \right\vert
= \left\vert \bfe_{i}^{\top} \bSigma \bH^{N+1} \bfe_{j} \right\vert
\leq \left\| \bfe_{i}^{\top} \bSigma \right\|_2
\left\| \bH^{N+1} \bfe_{j} \right\|_2
\leq \Tnorm{  \bSigma }_{2} \Tnorm{ \bH }_2^{N+1},
\eean
where $\bfe_{i} \in \{0,1\}^{D}$ denotes the standard basis vector whose $i$-th element is one and all other elements are zero.
Here, $\tnorm{\cdot}_2$ denotes the spectral norm  (or $\ell^2$ operator norm).
Since $\bx^{\top} \bSigma^{-1} \bx \geq (1-\epsilon KD) \| \bx \|_2^2$ for any $\bx \in \bbR^{D}$, the smallest eigenvalue of $\bSigma^{-1}$ is larger than $1-\epsilon  K D$; therefore, $ \tnorm{\bSigma}_2 \leq 1/(1-\epsilon KD)$.
Moreover, 
\bean
&& \Tnorm{\bH}_2 = \sup_{\| \bx \|_2 \leq 1} \| \bH \bx \|_2 = \sup_{\| \bx \|_2 \leq 1} \sqrt{ \sum_{i=1}^{D} \left( \sum_{j=1}^{D} H_{ij} x_j \right)^2 } \leq K \sup_{\| \bx \|_2 \leq 1} \sqrt{ \sum_{i=1}^{D} \left( \sum_{j=1}^{D} \vert x_j \vert \right)^2 }
\\
&& \leq D K \sup_{\| \bx \|_2 \leq 1} \| \bx \|_2 = DK.
\eean
Combining the last three displays, we have
\be
\left\vert \Sigma_{ij} - \sum_{k=0}^{N} \epsilon^{k} (\bH^{k})_{ij} \right\vert
\leq \frac{(\epsilon D K)^{N+1}}{1-\epsilon D K}
\label{eqlem1:errorbd}
\ee
for every $0 \leq \epsilon < (DK)^{-1}$ and $i,j \in [D]$.

Note that
\bean
(\bH^{k})_{ij} = \sum_{ \substack {i_0,\ldots,i_{k} \in [D] \\ i_0 = i, i_k = j }}  \prod_{l=0}^{k-1} H_{i_{l} i_{l+1} }
\eean
for $k \geq 1$, and $(\bH^{0})_{ij} = \delta_{ij}$, where $\delta_{ij}$ denotes the Kronecker delta.
For $i,j \in [D]$ with $i \neq j$ and $k \in \bbZ_{\geq 0}$ with $k < d_{G}(i,j)$, we have $(\bH^{k})_{ij} = 0$.
The proof proceeds by contradiction.
Suppose that $(\bH^k)_{ij} \neq 0$.
Then, there exists a sequence $(i^*_{0},\ldots,i^*_{k})$ such that $\prod_{l=0}^{k-1} H_{i^*_l i^*_{l+1}} \neq 0$ with $i_0^{*} = i$ and $i_{k}^{*} = j$.
This implies the existence of a path of length $k$ from vertex $i$ to $j$, which contradicts the assumption $k < d_{G}(i,j)$. 

For $i,j \in [D]$ with $d_{G}(i,j) < \infty$, we have
\bean
\sum_{k=0}^{d_{G}(i,j)} \epsilon^{k} (\bH^{k})_{ij} = \epsilon^{d_{G}(i,j)} \left( \bH^{d_{G}(i,j)} \right)_{ij}.
\eean
Combining with (\ref{eqlem1:errorbd}), we have
\bean
\left\vert \Sigma_{ij} -  \epsilon^{d_{G}(i,j)} \left( \bH^{d_{G}(i,j)} \right)_{ij}  \right\vert
\leq
 \frac{(\epsilon D K)^{d_{G}(i,j)+1}}{1-\epsilon D K}
\eean
for every $0 \leq \epsilon < (DK)^{-1}$ and $i,j \in [D]$ with $d_{G}(i,j) < \infty$.
Moreover, $(\bH^{k})_{ij} = 0$ for any $k \in \bbZ_{\geq 0}$ with $d_{G}(i,j) = \infty$.
Since $0 \leq \epsilon D K < 1$, combining with (\ref{eqlem1:errorbd}), we have $\Sigma_{ij} = 0$ for $d_{G}(i,j) = \infty$.

\end{proof}

\end{document}
